\section{Related work}
\label{sec:related}

LLaDA~\cite{llada2025} studies language modeling through masked
diffusion. DiffusionGemma~\cite{diffusiongemma} combines block-wise
diffusion with an encoder that caches context. Our contribution is an
empirical comparison of serving behavior for particular public models,
with explicit output-budget, latency-statistic, and quality constraints.
We do not claim that linear and staircase execution patterns are new,
or that these measurements isolate architecture from implementation.
PagedAttention and vLLM~\cite{vllm2023} provide the serving foundation;
using the same engine version removes one source of comparison bias,
but does not make its model-specific execution paths identical.

Grammar-constrained generation, including Outlines~\cite{outlines2024},
XGrammar~\cite{xgrammar2024}, and llguidance~\cite{llguidance},
addresses structured-output validity through token constraints.
This is complementary to choosing a decoding schedule. A completed
sampling run and an output that satisfies a grammar are different
properties; our unconstrained comparison cannot settle their joint
latency--quality tradeoff.

Speculative methods such as Medusa~\cite{medusa2024} and
EAGLE~\cite{eagle2024} verify multiple proposed tokens per iteration.
Their effective cost depends on acceptance, verification, and workload;
they need not obey the fitted slope of our standard AR baseline.
Our older n-gram prompt-lookup experiment on dense AR-3B is not a
comparison against those methods or a bound on optimized AR inference.
A tuned speculative AR-30B baseline is a useful next comparison.

Hoefler and Belli~\cite{hoefler2015} discuss experimental design and
reporting in systems benchmarking. The corrections in
Appendix~\ref{app:defects} provide a case study of how decoding
defaults, parsers, labels, and timing conventions affect a model-serving
comparison. They motivate explicit provenance and artifact checks,
without supporting an allegation about other authors' intentions.
